

\subsubsection{Freestanding Environment}

Relative command options: 



\subsubsection{Hosted Environment}

In Freestanding Environment, user can only use inline assembly, memory and arithmetic effect. However, the hosted includes at least interfaces provided by Operating System and Processor-registered mechanism, which provide much side-effect.

Usually to see demo `args':

\lstset{style=GlobalC}
\begin{lstlisting}[language=C, firstnumber=24]
#include <stdio.h>
#include <string.h>

int main(int argc, char *argv[], char* envp[]) {
	int i;
	printf("argc = %d;\n", argc);
	for (i = 0; i < argc; i++) {
		printf("arg[%d] = @\"%s\";\n", i, argv[i]);
	}
	if (argc > 1 && !strcmp("-e", argv[1])) for (i = 0; envp[i] != NULL; i++) {
		char* p = strchr(envp[i], '=');
		if (p) {
			*p = '\0';
			printf("env[%s] = @\"%s\";\n", envp[i], p + 1);
		}
		else printf("envp[%d] = @\"\"\"%s\"\"\";\n", i, envp[i]);
	}
	return 0;
}
\end{lstlisting}

But the ISO/IEC recommend the below forms for `main' function with external symbol \verb|_main|:

\verb|int main(int argc, char *argv[]) {/* ... */}| or 

\verb|int main(void) { }|

\textbf{EXAMPLE - Empty-return in Startup}

\begin{lstlisting}[language=C]
// ARINA RFW4 TEST on GCC Lin32
void main() { int a=2; }
// arina@aUbt:~$ echo $?
//:1
\end{lstlisting}

Learn the appliance with the utility `args'.

